\section*{\texorpdfstring{\S\CurrentExerciseGroup}{第{\CurrentExerciseGroup}节}}
\phantomsection
\addcontentsline{toc}{section}{第{\CurrentExerciseGroup}节习题}

\begin{enumerate}
\def\labelenumi{\arabic{enumi})}
\tightlist
\item
  设 H 是可积函数所成的一个集合，按关系 \(\leqslant\) 有向，这些函数 \(\geqslant 0\) 且满足 \(\sup N_{1}(f) < +\infty\)。设 g 是 H 的上包络；为使 g 可积且 \(f \in H\)
\end{enumerate}

在 \(L^{1}\) 中，g 的类 \(\widetilde{g}\) 是诸函数 \(f \in H\) 的类所成有向集的截口滤子之极限，这必须且只须：对每个 \(\varepsilon > 0\)，存在一个函数 \(f_{1} \in H\)、一个由下半连续可积函数组成且按 \(\leqslant\) 有向的集 B、以及一个映射 \(f \mapsto f^{*}\)，它定义在集 \(H_{1}\)（由诸函数 \(f \in H\) 中 \(\leqslant f_{1}\) 者组成）上且取值于集 B，使得 \(f \leqslant f^{*}\) 且 \(\mathrm{N}_{1}(f^{*} - f) \leqslant \varepsilon\) 对每个函数 \(f \in H_{1}\) 成立（利用 No.~4 的定理 3）。

\begin{enumerate}
\def\labelenumi{\arabic{enumi})}
\setcounter{enumi}{1}
\item
  设 \(\mu\) 是 \(\mathbf{R}\) 上的 Lebesgue 测度，\(\Omega\) 是在空间 \(\mathcal{L}^1\) 上赋予 \(\mathbf{R}\) 中逐点收敛拓扑所得的拓扑空间。试证：映射 \(f\mapsto \int fd\mu\)（从 \(\Omega\) 到 \(\mathbf{R}\)）在 \(\Omega\) 的任何点处都不连续。
\item
  设 I 是 \(\mathbf{R}\) 中的一个区间，\(\mathbf{f}\) 是 \(X \times I\) 到一个 Banach 空间 \(F\) 中的映射，满足：\(1^{\circ}\) 对每个 \(\alpha \in I\)，映射 \(t \mapsto \mathbf{f}(t, \alpha)\)（从 \(X\) 到 \(F\)）可积；\(2^{\circ}\) 对每个 \(t \in X\)，映射 \(\alpha \mapsto \mathbf{f}(t, \alpha)\) 有导数 \(\mathbf{f}_{\alpha}'(t, \alpha)\)（在 \(I\) 中）；\(3^{\circ}\) 存在一个可积函数 \(g \geqslant 0\)，使得 \(|\mathbf{f}_{\alpha}'(t, \alpha)| \leqslant g(t)\)（对一切 \(t \in X\) 与一切 \(\alpha \in I\)）。在这些条件下，试证：函数 \(\mathbf{u}(\alpha) = \int \mathbf{f}(t, \alpha) d\mu(t)\) 在 \(I\) 中可微，并且
\end{enumerate}

\[
\mathbf {u} ^ {\prime} (\alpha) = \int \mathbf {f} _ {\alpha} ^ {\prime} (t, \alpha) d \mu (t).
\]

\begin{enumerate}
\def\labelenumi{\arabic{enumi})}
\setcounter{enumi}{3}
\tightlist
\item
  设 \(\mu\) 是区间 \(X = [0,1]\) 上的 Lebesgue 测度。
\end{enumerate}

\begin{enumerate}
\def\labelenumi{\alph{enumi})}
\item
  在 X 中定义一个无处稠密的完全集 A，其测度为任意数 \(\alpha\)，满足 \(0 \leqslant \alpha < 1\)（利用 Cantor 三分集的构造方法）。
\item
  在 X 中定义一个由两两不交的无处稠密可积集组成的序列 \((\mathbf{A}_n)\)，使得 \(\mu (\mathrm{A}_n) = 2^{-n}\)，且使得 \(\mathbf{B}_n = \bigcup_{k = 1}^{n}\mathbf{A}_k\) 的每个余区间都含有 \(\mathrm{A}_{n + 1}\) 的一个测度 \(>0\) 的子集。若 \(\mathbf{A} = \bigcup_{n}\mathbf{A}_n\)，试证：\(\mathbf{A}\) 是测度为 1 的贫集，而 \(\mathbb{C}\mathbf{A}\) 是可忽略的非贫集。
\item
  设 \(\mathrm{H} = \bigcup_{n=0}^{\infty} \mathrm{A}_{2n+1}\)；试证：对每个开区间 \(\mathrm{I} \subset \mathrm{X}\)，I 与 H 以及与 CH 的交都有测度 \textgreater{} 0。设 f 是几乎处处等于特征函数 \(\varphi_{H}\) 的函数；试证：不存在任何由 X 上连续函数组成、在 X 的每点收敛且极限等于 f 的序列 \((f_{n})\)（注意 f 必然在 X 的每点不连续，并利用 GT, IX, §5 的习题 22）。
\end{enumerate}

\begin{enumerate}
\def\labelenumi{\arabic{enumi})}
\setcounter{enumi}{4}
\tightlist
\item
  设 \(\mu\) 是局部紧空间 \(X\) 上的一个正测度。对每个数值函数 \(f\)（无论有限与否、符号如何），只要它定义在 \(X\) 上，我们就用 \(\mu^{*}(f)\)（并称之为 \(f\) 的上积分）表示诸数 \(\mu(h)\)（对满足 \(h \geqslant f\) 的可积下半连续函数 h 而言）的下确界，当这样的函数存在时；在相反情形用 \(+\infty\)；当 \(f \geqslant 0\) 时，这个定义与 §1, No.~3 中的定义一致。我们用 \(\mu_{*}(f)\)（并称之为 \(f\) 的下积分）表示数 \(-\mu^{*}(-f)\)。
\end{enumerate}

\begin{enumerate}
\def\labelenumi{\alph{enumi})}
\item
  试证：若 \(f_{1}\) 与 \(f_{2}\) 是两个数值函数，使得几乎处处 \(f_{1}(x) \leqslant f_{2}(x)\)，则 \(\mu^{*}(f_{1}) \leqslant \mu^{*}(f_{2})\)。
\item
  设 \(f_{1}\) 与 \(f_{2}\) 是两个数值函数，使得 \(\mu^{*}(f_{1}) + \mu^{*}(f_{2})\) 有定义且 \(< +\infty\)；试证：\(f_{1}(x) + f_{2}(x)\) 几乎处处有定义，且 \(\mu^{*}(f_{1} + f_{2}) \leqslant \mu^{*}(f_{1}) + \mu^{*}(f_{2})\)（化归为在每点 \(f_{1}(x) < +\infty\) 且 \(f_{2}(x) < +\infty\) 的情形，借助 \(a\) )）。
\item
  若 \((f_n)\) 是一个递增的数值函数序列，使得从某个指标起 \(\mu^{*}(f_{n}) > -\infty\)，试证 \(\mu^{*}(\sup_{n} f_{n}) = \sup_{n} \mu^{*}(f_{n})\)。
\end{enumerate}

¶ 6) a) 设 \(f\) 是一个数值函数，使得 \(\mu^{*}(f)\)（习题 5）有限。试证：存在一个可积函数 \(f_{1} \geqslant f\)，使得 \(\mu(f_{1}) = \mu^{*}(f)\)；若 \(f_{2}\) 是另一个可积函数，满足 \(f_{2} \geqslant f\) 且 \(\mu(f_{2}) = \mu^{*}(f)\)，则 \(f_{1}\) 与 \(f_{2}\) 等价。

\begin{enumerate}
\def\labelenumi{\alph{enumi})}
\setcounter{enumi}{1}
\item
  要使数值函数 \(f\) 可积，必须且只须 \(\mu^{*}(f)\) 与 \(\mu_{*}(f)\) 有限且相等。
\item
  设 \(f\) 是一个数值函数，使得 \(\mu^{*}(f)\) 与 \(\mu_{*}(f)\) 都有限；设 \(g\) 与 \(h\) 是两个可积函数，满足 \(g \leqslant f \leqslant h\) 且 \(\mu(g) = \mu_{*}(f)\)、\(\mu(h) = \mu^{*}(f)\)。试证
\end{enumerate}

\[
\mu_ {*} (f - g) = \mu_ {*} (h - f) = 0 \quad \text { and } \quad \mu^ {*} (f - g) = \mu^ {*} (h - f) = \mu^ {*} (f) - \mu_ {*} (f).
\]

\begin{enumerate}
\def\labelenumi{\alph{enumi})}
\setcounter{enumi}{3}
\tightlist
\item
  设 \(f_{1}, f_{2}\) 是两个数值函数，使得数 \(\mu^{*}(f_{1}), \mu^{*}(f_{2})\)、\(\mu_{*}(f_{1})\) 与 \(\mu_{*}(f_{2})\) 都有限；试证
\end{enumerate}

\[
\mu_ {*} (f _ {1} + f _ {2}) \leqslant \mu_ {*} (f _ {1}) + \mu^ {*} (f _ {2}) \leqslant \mu^ {*} (f _ {1} + f _ {2})
\]

（若 \(g_2\) 是一个可积函数，满足 \(f_2 \leqslant g_2\) 且 \(\mu^*(f_2) = \mu(g_2)\)，则注意：对每个可积函数 \(h\)（满足 \(h \leqslant f_1 + f_2\)），有 \(h - g_2 \leqslant f_1\)）。由此推出：若 \(f_1\) 可积，则

\[
\mu^ {*} (f _ {1} + f _ {2}) = \mu (f _ {1}) + \mu^ {*} (f _ {2}), \quad \mu_ {*} (f _ {1} + f _ {2}) = \mu (f _ {1}) + \mu_ {*} (f _ {2}),
\]

且 \(\mu^{*}(f_{1} + f_{2}) = \mu^{*}\big(\sup (f_{1},f_{2})\big) + \mu^{*}\big(\inf (f_{1},f_{2})\big)\)（对后一关系式，利用 §1 的习题 1）。

\begin{enumerate}
\def\labelenumi{\alph{enumi})}
\setcounter{enumi}{4}
\tightlist
\item
  设 \(f\) 是一个可积函数。要使函数 \(g\)（\(\mu^{*}(g)\) 有限）可积，必须且只须 \(\mu(f) = \mu^{*}(g) + \mu^{*}(f - g)\)（若 \(g_{1}\) 是一个可积函数，满足 \(g \leqslant g_{1}\) 且 \(\mu^{*}(g) = \mu(g_{1})\)，注意 \(f - g_{1} \leqslant f - g\)）。
\end{enumerate}

¶ 7) 设 \(\mu\) 是 X 上的一个正测度。对每个子集 \(A \subset X\)，特征函数 \(\varphi_A\) 的下积分（习题 5）称为 A 的内测度，记作 \(\mu_*(A)\)。

\begin{enumerate}
\def\labelenumi{\alph{enumi})}
\item
  试证：\(\mu_{*}(\mathrm{A})\) 是含于 A 的紧集的测度之上确界（按定理 4 的方式论证）。
\item
  对 X 的每个外测度有限的子集 A，试证：存在两个可积集 \(A_{1}, A_{2}\)，使得 \(A_{1} \subset A \subset A_{2}\) 且 \(\mu_{*}(A) = \mu(A_{1})\)、\(\mu^{*}(A) = \mu(A_{2})\)。要使 A 可积，必须且只须 \(\mu^{*}(A)\) 与 \(\mu_{*}(A)\) 有限且相等。用同样的记号，试证
\end{enumerate}

\[
\mu_ {*} (\mathrm{A} \cap \mathbb {C} \mathrm{A} _ {1}) = \mu_ {*} (\mathrm{A} _ {2} \cap \mathbb {C} \mathrm{A}) = 0
\]

以及

\[
\mu^ {*} (\mathrm{A} \cap \mathbb {C} \mathrm{A} _ {1}) = \mu^ {*} (\mathrm{A} _ {2} \cap \mathbb {C} \mathrm{A}) = \mu^ {*} (\mathrm{A}) - \mu_ {*} (\mathrm{A}).
\]

\begin{enumerate}
\def\labelenumi{\alph{enumi})}
\setcounter{enumi}{2}
\item
  设 A 是一个可积集；试证：对每个集 \(\mathbf{B} \subset \mathbf{A}\)，有 \(\mu(\mathbf{A}) = \mu^{*}(\mathbf{B}) + \mu_{*}(\mathbf{A} \cap \mathbb{C}\mathbf{B})\)。
\item
  设 A 与 B 是两个不相交的外测度有限的集。试证：若 \(\mathbf{C} = \mathbf{A} \cup \mathbf{B}\)，则
\end{enumerate}

\[
\mu_ {*} (\mathrm{A}) + \mu_ {*} (\mathrm{B}) \leqslant \mu_ {*} (\mathrm{C}) \leqslant \mu_ {*} (\mathrm{A}) + \mu^ {*} (\mathrm{B}) \leqslant \mu^ {*} (\mathrm{C}) \leqslant \mu^ {*} (\mathrm{A}) + \mu^ {*} (\mathrm{B})
\]

以及

\[
\mu_ {*} (\mathrm{C}) - \mu_ {*} (\mathrm{A}) - \mu_ {*} (\mathrm{B}) \leqslant \mu^ {*} (\mathrm{A}) + \mu^ {*} (\mathrm{B}) - \mu^ {*} (\mathrm{C})
\]

（对后一不等式，化归为 \(\mu_{*}(\mathrm{A}) = \mu_{*}(\mathrm{B}) = 0\) 的情形，借助 \(b\) )）；若 \(\mathrm{A}_2\) 与 \(\mathrm{B}_2\) 是可积集，满足 \(\mathrm{A} \subset \mathrm{A}_2\)、\(\mathrm{B} \subset \mathrm{B}_2\)、\(\mu^{*}(\mathrm{A}) = \mu(\mathrm{A}_2)\)、\(\mu^{*}(\mathrm{B}) = \mu(\mathrm{B}_2)\)，试证 \(\mu_{*}(\mathrm{C}) \leqslant \mu(\mathrm{A}_2 \cap \mathrm{B}_2)\)。

¶ 8) 若把环面 T = R/Z 与 R 的区间 {[}0,1 {[} 典范地等同起来，该区间上的 Lebesgue 测度 μ 就是 T 上的一个测度；对每个集 A ⊂ T 及每个 z ∈ T，有 μ*(A + z) = μ*(A)。

\begin{enumerate}
\def\labelenumi{\alph{enumi})}
\item
  试证：存在子群 \(\mathrm{H}_0\)（\(\mathbf{T}\) 的），使得诸集 \(\mathrm{H}_n = r_n + \mathrm{H}_0\)（其中 \(r_n\) 取遍含于 {[}0,1{[} 的有理数所成之集）构成 \(\mathbf{T}\) 的一个分割（把 \(\mathbf{R}\) 看作 \(\mathbf{Q}\) 上的向量空间，注意在 \(\mathbf{R}\) 中子空间 \(\mathbf{Q}\) 有一个补子空间）。
\item
  设 H 是任意有限多个集 \(H_{n}\) 的并组成的集；试证：H 不可积且 \(\mu_{*}(H)=0\)（注意加群 Q/Z 中由有限个元素生成的每个子群在 Q/Z 中都有无限指数；由此推出，存在 T 的一个分成可数无穷多个集 \(P_{n}\) 的分割，其中每个集都含有形如 \(z+H\) 的集）。
\item
  由 \(b\) ) 出发，给出一个递减序列 \((\mathrm{A}_n)\)（\(\mathbf{T}\) 的子集的序列）的例子，其交为空集，且 \(\mu^{*}(\mathrm{A}_{n}) = 1\)（对所有 \(n\)）。
\item
  设 A 是一个可积集，满足 \(H_{0} \subset A\) 且 \(\mu(A) = \mu^{*}(H_{0}) > 0\)；试证：对每个可积集 \(B \subset A\)，集 \(B_{1} = B \cap H_{0}\) 与 \(B_{2} = B \cap C H_{0}\) 构成 B 的一个分割，使得 \(\mu^{*}(B_{1}) = \mu^{*}(B_{2}) = \mu(B)\) 且 \(\mu_{*}(B_{1}) = \mu_{*}(B_{2}) = 0\)（利用习题 7 b)）。
\end{enumerate}

¶ 9) a) 设 \(\mu\) 是局部紧空间 \(X\) 上的一个正测度。在 Banach 空间 \(\widetilde{\mathcal{F}}^1\)（由 \(\mathcal{F}^1\) 中函数的等价类组成）中，设 \(G\) 是一个包含子空间 \(L^{1}\) 的闭线性子空间；设 G 是由类属于 G 的函数组成的 \(F^{1}\) 的闭线性子空间。试证：可把线性形式 \(\mu(f)\) 延拓到 G 上，使不等式 \(|\mu(f)| \leqslant N_{1}(f)\) 仍成立（利用 Hahn--Banach 定理）。试证：对这样延拓的积分，§3 的定理 3 仍然有效。

\begin{enumerate}
\def\labelenumi{\alph{enumi})}
\setcounter{enumi}{1}
\tightlist
\item
  假设 G 是 \(\mathbf{L}^1\) 与一个有限维子空间 H 的直和。试证：在 \(\mathcal{G}\) 中，Lebesgue 定理（定理 2）仍然有效。（设 \((f_n)\) 是 \(\mathcal{G}\) 中函数的一个序列，几乎处处趋于 \(f\)，且满足 \(|f_n| \leqslant g\)，其中 \(g \geqslant 0\) 且 \(\mathrm{N}_1(g) < +\infty\)。设 \((\widetilde{u}_k)_{1 \leqslant k \leqslant m}\) 是 H 的一个基，并设 \(f_n = \sum_{k=1}^{m} \alpha_{nk} u_k + h_n\)，其中 \(h_n \in \mathcal{L}^1\)。利用 \(G\) 是 \(L^1\) 与 \(H\) 的拓扑直和这一事实，试证诸 \(\alpha_{nk}\) 一致有界；必要时过渡到 \((f_n)\) 的一个子序列，化归为每个序列 \((\alpha_{nk})_{n \geqslant 1}\) 都有极限的情形；由此推出 \(h_n(x)\) 这时几乎处处趋于一个极限，并对序列 \((h_n)\) 应用 Lebesgue 定理；最后注意 \((\widetilde{f}_n)\) 的两个子序列在 \(G\) 中不可能趋于不同的极限。
\end{enumerate}

¶ 10) a) 设 E 是一个 Riesz 空间，μ 是 E 上的一个正线性形式，使得关系 μ(\textbar x\textbar) = 0 蕴含 x = 0；于是 μ(\textbar x\textbar) 是 E 上的一个范数，我们假设 E 赋此范数后是完备的。由此推出 E 是完全格序的（Ch. II, §2, 习题 8 e)）。从而（Ch. II, §1, 习题 12 与 13）存在一个局部紧空间 X，它是一族 (Kα) 紧 Stone 空间之和，使得 E 同构于 X 上的一个包含空间 K(X) 的连续数值函数（无论有限与否）所成的空间。把 E 与这个空间等同起来，μ 在 K(X) 上的限制就是 X 上的一个正测度。试证：E 典范地同构于空间 L¹(μ)；更精确地说，对每个定义在 X 上的 μ-可积函数 g，存在一个且仅一个函数 f ∈ E，它关于测度 μ 与 g 等价（注意 E 的每个元素 ≥ 0 都是 K(X) 的元素的一个递增序列的上确界，且 K(X) 在 L¹(μ) 中稠密）。

\begin{enumerate}
\def\labelenumi{\alph{enumi})}
\setcounter{enumi}{1}
\tightlist
\item
  由 a) 推出：对每个紧空间 K，存在一个局部紧空间 S（一族紧 Stone 空间的拓扑和）及 S 上的一个正测度 \(\nu\)（其支集等于 S），使得 K 上的测度所成的完全格序空间 \(\mathcal{M}(\mathrm{K})\)（赋以范数 \(\|\mu\|\)）同构于 \(\mathcal{L}^{1}(\nu)\)（考虑线性形式 \(\mu\mapsto\mu(\mathrm{K})\)，在 \(\mathcal{M}(\mathrm{K})\) 上）。
\end{enumerate}

\begin{enumerate}
\def\labelenumi{\arabic{enumi})}
\setcounter{enumi}{10}
\tightlist
\item
  \begin{enumerate}
  \def\labelenumii{\alph{enumii})}
  \tightlist
  \item
    设 \(\Gamma\) 是一个集 \(A\) 的子集的任意一个集合。设 \(\Psi\) 是 \(A\) 的形如下式的子集所成之集
  \end{enumerate}
\end{enumerate}

\[
\mathrm{X} _ {1} \cap \mathrm{X} _ {2} \cap \dots \cap \mathrm{X} _ {m} \cap \mathbb {C} \mathrm{X} _ {m + 1} \cap \dots \cap \mathbb {C} \mathrm{X} _ {m + p},
\]

其中诸 \(X_{i}\) 是 \(\Gamma\) 中的集，m 是任意整数 \(\geqslant 1\)，p 是任意整数 \(\geqslant 0\)。试证：最小集环 \(\Phi\)（包含 \(\Gamma\)）是 \(\Psi\) 中之集的有限并所成之集。

\begin{enumerate}
\def\labelenumi{\alph{enumi})}
\setcounter{enumi}{1}
\item
  设 \(\Delta\) 是 \(\Gamma\) 中之集的有限交所成之集。试证：对每个向量空间 F，由 \(\Gamma\) 生成的集环中之集的特征函数的线性组合（系数取自 F）所成之集，等同于 \(\Delta\) 中之集的特征函数的线性组合所成之集。
\item
  设 X 是拓扑空间，\(\Gamma\) 是 X 的紧子集所成之集。试证：由 \(\Gamma\) 生成的集环等同于 X 的形如 \(A \cap CB\) 的子集的有限并所成之集，其中 A 与 B 是紧集。
\end{enumerate}

\begin{enumerate}
\def\labelenumi{\arabic{enumi})}
\setcounter{enumi}{11}
\tightlist
\item
  设 \(X\) 是一个 Hausdorff 拓扑空间。对每个偶 \((K, U)\)（由一个紧集 \(K\) 与一个开集 \(U\)——\(X\) 中的——组成），设 \(I(K, U)\) 是子集 \(M \subset X\)（满足 \(K \subset M \subset U\)）所成之集，并设 \(\mathcal{T}\) 是由 \(\mathfrak{P}(X)\) 的诸子集 \(I(K, U)\) 在 \(\mathfrak{P}(X)\) 上生成的拓扑。
\end{enumerate}

\begin{enumerate}
\def\labelenumi{\alph{enumi})}
\item
  试证：每个集 I(K, U) 在 \(\mathfrak{P}(X)\) 中既开又闭；由此推出，\(\mathfrak{P}(X)\) 赋以拓扑 \(\mathcal{T}\) 后是一个完全正则且全不连通的空间。
\item
  要使映射 \(M \mapsto CM\)（从 \(\mathfrak{P}(X)\) 到其自身）连续（关于拓扑 T），必须且只须 X 是紧的。
\item
  取 X 为 R 的区间 {[}0,1{]}。试证：映射 \((M,N)\mapsto M\cup N\) 与 \((M,N)\mapsto M\cap N\)（从 \(\mathfrak{P}(X)\times\mathfrak{P}(X)\) 到 \(\mathfrak{P}(X)\) 中）关于拓扑 T 不连续。
\item
  假设 X 是局部紧的。试证：T 在 X 的紧子集所成之集上诱导的拓扑，细于按 GT, II, §1 习题 5 的程序由 X 上一个一致结构导出的拓扑；除非 X 是离散空间，这两个拓扑不可能相同。
\end{enumerate}

¶ 13) a) 设 X 是局部紧空间，\(\Gamma\) 是由相对紧的集组成的 X 拓扑的一个基。设 \(\mathcal{U}\) 是与 X 的拓扑相容的一个一致结构，\(\mathfrak{S}\) 是该结构的一个邻域基本系。设 \(M \mapsto \lambda(M)\) 是一个有限数值函数 \(\geqslant 0\)（定义在 \(\Gamma\) 上）。对每个紧集 \(K \subset X\) 与每个邻域 \(V \in \mathfrak{S}\)，设 \(\alpha_V(K)\) 是数 \(\sum_{i} \lambda(U_i)\)

对 K 的所有有限覆盖 \((\mathrm{U}_{i})\)（由 \(\Gamma\) 中关于 V 为细小的集组成）；假设当 V 取遍 S 时，\(\alpha(\mathrm{K})\)（诸数 \(\alpha_{\mathrm{V}}(\mathrm{K})\) 的上确界）有限。试证：存在 X 上的一个测度 \(\mu\)（且仅一个），使得对每个紧集 K 有 \(\mu(\mathrm{K}) = \alpha(\mathrm{K})\)（利用定理 5）。

\begin{enumerate}
\def\labelenumi{\alph{enumi})}
\setcounter{enumi}{1}
\tightlist
\item
  设 \(\Psi\) 是 \(X\) 的 Borel 子集所成之集，\(\alpha\) 是 \(\Psi\) 到 \([0, +\infty]\) 中的一个映射，满足下列条件：
\end{enumerate}

\begin{enumerate}
\def\labelenumi{(\roman{enumi})}
\item
  若 \(\mathrm{B}_1, \mathrm{B}_2\) 是 \(X\) 的两个不交的 Borel 子集，则 \(\alpha(\mathrm{B}_1 \cup \mathrm{B}_2) = \alpha(\mathrm{B}_1) + \alpha(\mathrm{B}_2)\)。\\
\item
  若 \(B\) 是 \(X\) 的一个紧子集，则 \(\alpha(B) < +\infty\)。
\item
  若 \(\mathbf{B}\) 是 \(\mathbf{X}\) 的一个 Borel 子集，则 \(\alpha(\mathbf{B}) = \inf \alpha(\mathbf{U})\)，其中 \(\mathbf{U}\) 取遍 \(\mathbf{X}\) 中包含 \(\mathbf{B}\) 的开子集所成之集。
\end{enumerate}

那么，存在一个且仅一个 X 上的正测度 \(\mu\)，使得对 X 的每个可被一列紧集覆盖的 Borel 子集 B 有 \(\alpha(B) = \mu^{*}(B)\)。c) 对每个 \(B \in \Phi\)，若 B 可被一列紧集覆盖则令 \(\alpha(B) = 0\)，在相反情形令 \(\alpha(B) = +\infty\)。这时 \(b\) 的条件 (i)、(ii)、(iii) 满足。我们有 \(\mu = 0\)，因此若 B 不能被一列紧集覆盖，则 \(\alpha(B) \neq \mu^{*}(B)\)。

\begin{enumerate}
\def\labelenumi{\arabic{enumi})}
\setcounter{enumi}{13}
\item
  设 \((\mathbf{A}_n)\) 是一个可积集序列，使得 \(\sum_{n} \mu(\mathbf{A}_n) < +\infty\)。对每个整数 \(k\)，设 \(\mathbf{G}_k\) 是由 \(x \in X\) 组成的集，其中 \(x \in \mathbf{A}_n\) 至少对 \(k\) 个值 \(n\) 成立；试证：\(\mathbf{G}_k\) 可积，且 \(k \cdot \mu(\mathbf{G}_k) \leqslant \sum_{n} \mu(\mathbf{A}_n)\)。
\item
  设 \(\mu\) 是局部紧空间 \(X\) 上的一个有界正测度，\((A_n)\) 是 \(X\) 中可积集的一个序列，使得 \(\inf_{n} \mu(A_n) = m > 0\)。试证：集 \(B\)——\(X\) 中属于无穷多个集 \(A_n\) 的点所成之集——是可积的，且 \(\mu(B) \geqslant m\)。
\end{enumerate}

¶ 16) 设 X 是一个完全正则空间，\(\mathcal{C}(X)\)（相应地，\(\mathcal{C}^b(X)\)）是 X 上连续（相应地，连续且有界）数值函数的 Riesz 空间。

\begin{enumerate}
\def\labelenumi{\alph{enumi})}
\item
  试证：若线性形式 \(\lambda\)（\(\mathcal{C}(\mathrm{X})\) 上的）关于紧收敛拓扑连续，则它是相对有界的。
\item
  设 \(\lambda\) 是 \(\mathcal{C}(\mathrm{X})\) 上的一个正线性形式。试证：若 \(\lambda\) 在 \(\mathcal{C}^b (\mathrm{X})\) 上为零，则它在 \(\mathcal{C}(\mathrm{X})\) 上为零（设 \(\varphi\) 是 \(\mathcal{C}(\mathrm{X})\) 到商 Riesz 空间 \(\mathcal{C}(\mathrm{X}) / \mathcal{C}^b (\mathrm{X})\)（Ch. II, §1, 习题 4）上的典范映射；试证：若 \(h\) 是一个连续函数 \(\geqslant 0\) 且在 \(\mathrm{X}\) 上不有界，则 \(n\varphi (h)\leqslant \varphi (h^{2})\)（对每个整数 \(n > 0\)））。
\item
  设 \(\beta X\) 是 \(X\) 的 Stone--Čech 紧化，即通过在 \(X\) 上赋予以使 \(\mathcal{C}^b(X)\) 中函数都一致连续的最粗一致结构、再把如此得到的一致空间完备化而得的紧空间；这时每个函数 \(f \in \mathcal{C}(X)\) 都可连续延拓为连续函数 \(\widetilde{f}\)（在 \(\beta X\) 上，可能取无穷值）（考虑 \(f/(1 + |f|)\)）。若 \(\lambda\) 是 \(\mathcal{C}(X)\) 上的一个正线性形式，则 \(\lambda\) 在 \(\mathcal{C}^b (\mathrm{X})\) 上的限制形如 \(f\mapsto \mu (\widetilde{f})\)，其中 \(\mu\) 是 \(\beta X\) 上的一个正测度；试证：对每个函数 \(f\in \mathcal{C}(\mathrm{X})\)，\(\widetilde{f}\) 关于 \(\mu\) 可积，且 \(\lambda (f) = \mu (\widetilde{f})\)（利用 \(b\) )，并注意每个函数 \(\geqslant 0\)（属于 \(\mathcal{C}(\mathrm{X})\) 的）都是 \(\mathcal{C}^b (\mathrm{X})\) 中函数的一个序列的上包络）。
\item
  试证：每个点 \(x_{0}\)（属于 \(\mu\) 的支集而不属于 X 的）都具有下列性质：对每个递减序列 \((V_{n})\)（由 \(x_{0}\) 的邻域组成），诸 \(V_{n}\) 的交至少含有 X 的一个点。逆命题成立。
\item
  由 \(d\) ) 推出：若 \(X\) 局部紧且无穷远处可数，则 \(\mu\) 的支集含于 \(X\)（与 \(h\) ) 比较）。
\item
  若 X 是离散的，试证：\(\mu\) 的支集有限（在相反情形，造一个定义在 X 上的函数 \(f \geqslant 0\)，使 \(\widetilde{f}\) 不是 \(\mu\) -可积的）。
\item
  试证：若线性形式 \(\lambda\)（\(\mathcal{C}(\mathbf{X})\) 上的）是正的且关于紧收敛拓扑连续，则 \(\mu\) 的支集含于 X（参见 Ch. III, §2, No.~3, 命题 11）。
\item
  设 \(X_{0}\) 是把 GT, I, §9, 习题 11 中定义的局部紧空间 X 添加一个无穷远点 \(\omega\) 所得的紧空间，Y 是 \(\overline{R}\) 中由整数 \(\geqslant 0\) 与 \(+\infty\) 组成的子空间，Z 是点 \((\omega, +\infty)\) 在乘积空间 \(X_{0} \times Y\) 中的余空间（局部紧）。试证：Z 上每个连续数值函数都有界，且 \(f(z)\) 当 z 趋于 Z 的无穷远点 \((\omega, +\infty)\) 时趋于一个极限。由此推出：若 \(\mu\) 是 Z 上由把质量 \(1/2^{n}\) 放在每个点 \((\omega, n)\) 处（对每个 \(n \geqslant 0\)）所定义的测度，则 Z 上每个连续函数都 \(\mu\) -可积，但 \(\mu\) 的支集不是紧的。
\end{enumerate}

¶ 17) 设 X 是 GT, I, §9, 习题 11 中定义的局部紧空间。

\begin{enumerate}
\def\labelenumi{\alph{enumi})}
\tightlist
\item
  设 H 是赋以紧收敛拓扑的局部凸空间 \(\mathcal{C}(\mathrm{X};\mathbf{C})\) 的一个紧子集；试证：诸函数 \(f \in H\) 在 X 上一致有界，且存在 \(c \in X\)，使得对每个 \(x \geqslant c\)，H 中所有函数都取常值。（用反证法论证；若 H 中函数在 X 上不一致有界，就会存在 X 的点的一个递增序列 \((x_{n})\)，使得 \(\sup(|\mathrm{H}(x_{n})|) \geqslant n\)，其中 \(|\mathrm{H}(x_{n})|\) 表示诸数 \(|f(x_{n})|\)（对 \(f \in H\)）所成之集；注意序列 \((x_{n})\) 在 X 中收敛。类似地，对每个 \(x \in X\)，设 \(\delta(x)\) 是所有函数 \(f \in H\) 在区间 \([x, \rightarrow[; \delta(x)\) 上的振幅的上确界；它有限且递减，因此存在 \(d \in X\)，使得 \(\delta(x)\) 对 \(x \geqslant d\) 取常值。为证这个常数 \(\beta\) 为零，像前面那样用反证法，造出 X 的点的两个递增序列 \((s_{n}), (t_{n})\)，满足 \(s_{n} \leqslant t_{n} \leqslant s_{n+1} \leqslant t_{n+1}\)，以及 H 中函数的一个序列 \((f_{n})\)，满足 \(|f_{n}(s_{n}) - f_{n}(t_{n})| \geqslant \beta/2\)。
\end{enumerate}

特别地，\(\mathcal{C}(\mathrm{X};\mathbf{C})\) 是 \(\mathcal{K}(\mathrm{X};\mathbf{C})\) 与 \(C\cdot1\) 的直和。

\begin{enumerate}
\def\labelenumi{\alph{enumi})}
\setcounter{enumi}{1}
\item
  试证：在 X 上每个测度都有紧支集（对每个 \(x \in X\)，设 \(f_x\) 是区间 \(\leftarrow, x\) {]} 的特征函数，它连续且具有紧支集；考虑 X 上的递增函数 \(|\mu|(f_x)\)）。
\item
  由 a) 与 b) 推出：在 \(\mathcal{M}(\mathrm{X};\mathbf{C}) = \mathcal{M}^{1}(\mathrm{X};\mathbf{C}) = \mathcal{C}'(\mathrm{X};\mathbf{C})\) 中，关于 \(\mathcal{T}\)（\(\mathcal{C}(\mathrm{X};\mathbf{C})\) 的紧子集上一致收敛拓扑）的有界集，就是关于超强拓扑（Ch. III, §1, 习题 15）的有界集。试证：\(\mathcal{C}'(\mathrm{X};\mathbf{C})\) 关于拓扑 \(\mathcal{T}\) 不是拟完备的（考虑测度 \(\varepsilon_{x}\)，对 \(x\in \mathrm{X}\)）；由此推出，关于紧收敛拓扑，\(\mathcal{C}(\mathrm{X};\mathbf{C})\) 既非桶式也非有界型（参见 TVS, III, §4, No.~2, 定理 1 的推论 4 及下面的习题 18）。
\item
  试证：在 \(\mathcal{K}(\mathrm{X};\mathbf{C})\) 上，拓扑 \(\mathcal{T}_0\)（诸 Banach 空间 \(\mathcal{K}(\mathrm{X},\mathrm{K};\mathbf{C})\) 的拓扑的直接极限）等同于一致收敛拓扑，且赋此拓扑后 \(\mathcal{K}(\mathrm{X};\mathbf{C})\) 是完备的。（设 V 是 \(\mathcal{T}_0\) 的一个 0 的邻域；对每个 \(x\in X\)，设 \(r_x\) 是最大的数 \(>0\)，使得 V 包含每个连续函数 \(f\)，其支集含于 \([ \leftarrow ,x]\) 且满足 \(\| f\| \leqslant r_x\)。试证 \(r_x\) 在 X 中的下确界 \(>0\)，为此证明不可能存在 X 的点的递增序列 \((x_n)\)，使得 \(\lim_{n\to \infty}r_{x_n} = 0\)。）
\end{enumerate}

\begin{enumerate}
\def\labelenumi{\arabic{enumi})}
\setcounter{enumi}{17}
\tightlist
\item
  设 \(X\) 是一个仿紧的局部紧空间。
\end{enumerate}

\begin{enumerate}
\def\labelenumi{\alph{enumi})}
\item
  试证：空间 \(\mathcal{C}(\mathrm{X};\mathbf{C})\) 同构于一族 Fréchet 空间的乘积，从而是桶式的。
\item
  要使 \(\mathcal{C}^{\prime}(X; \mathbf{C})\) 的子集 H 关于紧收敛拓扑有界，必须且只须存在 X 的一个紧子集 K，使得 \(\operatorname{Supp}(\mu) \subset K\)（对一切 \(\mu \in H\)），且 H 关于超强拓扑有界（与习题 17 c) 比较）。这时紧收敛拓扑在 H 上诱导的拓扑等同于淡拓扑。
\item
  试证：在集 \(\mathcal{M}_{+}(\mathrm{X})\cap \mathcal{C}'(\mathrm{X};\mathbf{C})\) 上，由 \(\mathcal{C}'(\mathrm{X};\mathbf{C})\) 上紧收敛拓扑诱导的拓扑等同于淡拓扑。当 \(X\) 不是仿紧时（习题 17），结论还成立吗？
\end{enumerate}

\begin{enumerate}
\def\labelenumi{\arabic{enumi})}
\setcounter{enumi}{18}
\tightlist
\item
  在 \(\mathcal{C}(\mathrm{X};\mathbf{C})\) 上赋以紧收敛拓扑，在 \(\mathcal{C}'(\mathrm{X};\mathbf{C})\) 上赋以 \(\mathcal{C}(\mathrm{X};\mathbf{C})\) 的紧子集上一致收敛拓扑。设 \(\mathfrak{S}\) 是 \(\mathcal{C}(\mathrm{X};\mathbf{C})\) 的紧子集所成之集；试证：双线性映射 \((g,\mu)\mapsto g\cdot \mu\)（从 \(\mathcal{C}(\mathrm{X};\mathbf{C})\times \mathcal{C}'(\mathrm{X};\mathbf{C})\) 到 \(\mathcal{C}'(\mathrm{X};\mathbf{C})\) 中）是 \(\mathfrak{S}\) -亚连续的。设 \(\mathcal{T}\) 是 \(\mathcal{C}'(\mathrm{X};\mathbf{C})\) 的有界子集所成之集；试证：若 X 是仿紧的，则前面的双线性映射也是 \(\mathcal{T}\) -亚连续的；当不再假设 X 是仿紧时（习题 17），后一性质还成立吗？
\end{enumerate}

¶ 20) 设 X, Y 是两个仿紧的局部紧空间。在空间 \(\mathcal{C}'(X; \mathbf{C})\) 与 \(\mathcal{C}'(Y; \mathbf{C})\) 上赋以紧收敛拓扑。

\begin{enumerate}
\def\labelenumi{\alph{enumi})}
\item
  试证：当在 \(\mathcal{C}^{\prime}(X\times Y;\mathbf{C})\) 上赋以紧收敛拓扑时，映射 \((\mu,\nu)\mapsto\mu\otimes\nu\) 是 \((\mathfrak{S},\mathfrak{T})\) -亚连续的，其中 S 表示 \(\mathcal{C}^{\prime}(X;\mathbf{C})\) 的有界子集所成之集，T 表示 \(\mathcal{C}^{\prime}(Y;\mathbf{C})\) 的有界子集所成之集（利用习题 18 b)）。
\item
  试证：当 X 与 Y 无穷远处可数且在 \(\mathcal{C}^{\prime}(X\times Y;C)\) 上赋以紧收敛拓扑时，映射 \((\mu,\nu)\mapsto\mu\otimes\nu\) 连续。（受 Ch. III, §4 习题 3 的证明启发，在 X（相应地 Y）中引入从属于一个局部有限开覆盖的连续单位分解，按 TVS, II, §4, 习题 9 a) 的方式论证。）
\item
  若 X 是离散的，空间 \(\mathcal{C}^{\prime}(X;\mathbf{C})\) 可等同于直和空间 \(\mathbf{C}^{(X)}\)（赋以最细局部凸拓扑）。由此给出一个例子：X 与 Y 都是离散的，X 不可数，且映射 \((\mu,\nu)\mapsto\mu\otimes\nu\) 当 \(\mathcal{C}^{\prime}(X\times Y;\mathbf{C})\) 赋以淡拓扑时不连续（参见 TVS, II, §4, 习题 9 b)）。
\end{enumerate}

\begin{enumerate}
\def\labelenumi{\arabic{enumi})}
\setcounter{enumi}{20}
\item
  设 \(\mathbf{X}\) 是一个局部紧空间，\(f\) 是一个数值函数 \(\geqslant 0\)（定义在 \(\mathbf{X}\) 上）。要使 \(\mu \mapsto \mu^{*}(f)\) 从 \(\mathcal{M}_{+}(\mathbf{X}) \cap \mathcal{C}'(\mathbf{X};\mathbf{C})\) 到 \(\overline{\mathbf{R}}\) 中的映射关于紧收敛拓扑连续，必须且只须 \(f\) 在 \(\mathbf{X}\) 上连续。
\item
  设 X, Y 是两个局部紧空间，\(\mu\) 是 X 上的一个有界测度，\(\nu\) 是 Y 上的一个有界测度，于是 \(\mu \otimes \nu\) 是 \(X \times Y\) 上的一个有界测度。试证：对每个函数 \(f \in \mathcal{C}^{0}(X \times Y; \mathbf{C})\)，函数 \(x \mapsto \int f(x, y) d\nu(y)\) 属于 \(\mathcal{C}^{0}(X; \mathbf{C})\)，并且
\end{enumerate}

\[
\int f (x, y) d \mu (x) d \nu (y) = \int d \mu (x) \int f (x, y) d \nu (y).
\]

¶ 23) 在空间 \(\mathbf{R}^n\) 上，设 \(\mu\) 是 Lebesgue 测度，\(|\mathbf{x}|\) 是一个范数，使该范数的单位球测度等于 1，\((\mathbf{x}_k)_{k \geqslant 1}\) 是一个有界可积集 B 中互异点的一个无穷序列，其中 \(\mu(B) = 1\)。对每个整数 \(m\)，用 \(d_m\) 表示数 \(|\mathbf{x}_i - \mathbf{x}_j|\)（对 \(1 \leqslant i < j \leqslant m\)）的下确界。试证

\(\liminf_{m \to \infty} md_m^n \leqslant \alpha_n^{-1}\)，其中

\[
\alpha_ {n} = 1 + n \int_ {0} ^ {1} \frac {(1 - t) ^ {n}}{1 + t} d t.
\]

（用反证法论证：假设对某个 \(\varepsilon > 0\)，存在一个 \(m_0\)，使得 \(md_m^n > h_n\)（对 \(m \geqslant m_0\)），其中 \(h_n = \alpha_n^{-1} + \varepsilon\)。对 \(1 \leqslant i < m\)，设 \(B_i\) 是以 \(\mathbf{x}_i\) 为中心、以 \(\frac{1}{2} h_n m^{-1/n}\) 为半径的球；对 \(m \leqslant i \leqslant 2^n m\)，设 \(B_i\) 是以 \(\mathbf{x}_i\) 为中心、以 \(\frac{1}{2} h_n (2i^{-1/n} - m^{-1/n})\) 为半径的球。试证：这 \(2^n m\) 个球 \(B_i\) 两两不交，并利用 Euler--Maclaurin 公式计算它们的并的测度。）
% semantic-fragments: legacy-input-seam
\relax{}
\setcounter{exercisegroup}{5}
\edef\CurrentExerciseGroup{\arabic{exercisegroup}}
